% --- ARCHIVO: Anexo_Prompts/prompts.tex ---

\chapter{Registro de \textit{Prompts} y Directrices de IA}
\label{app:prompts}

El presente anexo documenta de forma transparente la dimensión operativa del
flujo de trabajo descrito en el capítulo~\ref{cap:uso_ia}
(p.~\pageref{cap:uso_ia}). Se recogen, en orden cronológico de incorporación
al proceso, cinco \textit{prompts} representativos de las distintas
funciones desempeñadas por la asistencia LLM a lo largo del trabajo: el
\textit{prompt} maestro de configuración que fijó las reglas
ortotipográficas y el encuadre causal del colapso, dos \textit{prompts} de
estructuración documental ---reconciliación cronológica entre informes y
extracción normativa con anclaje a fuente---, y dos \textit{prompts}
correctores emitidos en el momento exacto en que el modelo incurrió en
inferencias incompatibles con la física del caso (Paradoja del UFLS y
\textit{Tap-Lag}). Los textos se reproducen con la misma redacción
restrictiva con la que fueron utilizados, conservando el español técnico y
las marcas léxicas que delimitaban el espacio de respuesta admisible para
el modelo. Cada \textit{prompt} se acompaña de una breve nota explicativa
que sitúa su función dentro del flujo de validación física documentado en
el capítulo metodológico.

\section*{\textit{Prompt}~1. Configuración maestra: encuadre causal y
reglas ortotipográficas LaTeX}
\label{prompt:maestro}

Este \textit{prompt} se introdujo como mensaje de sistema al inicio de
cada sesión de trabajo, fijando dos perímetros que el modelo debía
respetar a lo largo de toda la conversación: la tesis causal del colapso
---inestabilidad de tensión capacitiva, no déficit primario de inercia
rotacional--- y las convenciones tipográficas LaTeX del documento. Su
función fue prevenir, antes incluso de cualquier consulta sustantiva, las
dos categorías de error más frecuentes detectadas en interacciones
preliminares: la deriva narrativa hacia el paradigma frecuencia-potencia
activa y la inserción sistemática de unidades dentro del modo matemático.

\begin{quote}
    \itshape\small
    ``Vas a actuar como asistente de redacción técnica para un Trabajo de
    Fin de Grado en Ingeniería de la Energía sobre el apagón del sistema
    eléctrico ibérico del 28 de abril de 2025. Antes de aceptar cualquier
    consulta posterior, fija los siguientes dos perímetros y respétalos
    sin excepción durante toda la sesión.

    \medskip
    \textbf{(A) Encuadre causal del colapso (no negociable).} El cero de
    tensión del 28A no fue un colapso primario de frecuencia por déficit
    de inercia rotacional. Fue una inestabilidad de tensión en régimen
    capacitivo, precipitada por la maniobra de mallado del Operador del
    Sistema en condiciones de baja fortaleza síncrona, agravada por las
    limitaciones de observabilidad sobre la red de 220~kV (fenómeno
    Tap-Lag). La inercia zonal baja (1,3~s en el área Sur, 1,84~s en el
    Centro) es condición habilitante, no causa raíz. Si en algún momento
    tu razonamiento se desvía hacia la lógica frecuencia-potencia activa
    propia de los apagones clásicos, detente, declara la desviación y
    espera reformulación del prompt. No completes razonamientos
    causales que contradigan este encuadre.

    \medskip
    \textbf{(B) Reglas ortotipográficas LaTeX (estrictas).} Las
    magnitudes con unidades van SIEMPRE fuera del modo matemático.
    Correcto: 400~kV, 82~\%, 2,3~s, 1.050~MVAr, 12:33~CEST. Incorrecto:
    \$400\textbackslash text\{ kV\}\$, \$82\textbackslash,\textbackslash\%\$.
    Solo las variables propiamente dichas van dentro de \$\dots\$:
    \$H\$, \$V\$, \$P\$, \$Q\$, \$f\$, \$\textbackslash delta\$. Los
    decimales se escriben con coma, no con punto. Las citas se agrupan
    siempre en un único corchete: \textbackslash cite\{a, b\}, nunca
    \textbackslash cite\{a\}\textbackslash cite\{b\}. Las rutas de
    \textbackslash includegraphics, \textbackslash label y
    \textbackslash ref no se modifican bajo ninguna circunstancia.

    \medskip
    Si una consulta posterior contradice (A) o (B), prevalece la regla.
    Si detectas ambigüedad, formula una pregunta de clarificación antes
    de generar texto.''
\end{quote}

\section*{\textit{Prompt}~2. Reconciliación cronológica de
\textit{timestamps} entre informes primarios}
\label{prompt:cronologia}

Los cuatro informes primarios sobre el incidente ---Consejo de Seguridad
Nacional, Red Eléctrica de España, IIT-ICAI/AELEC y ENTSO-E--- describen
los mismos hechos físicos con distintos niveles de granularidad temporal
y, ocasionalmente, con desfases internos en la franja de los 100--200~ms.
Este \textit{prompt} se utilizó para construir la tabla maestra de
\textit{timestamps} unificados que articula la cronología expuesta en el
capítulo~\ref{cap:analisis_incidente}.

\begin{quote}
    \itshape\small
    ``Tienes acceso a los cuatro informes primarios del 28A en formato
    PDF. Vas a construir una cronología unificada del intervalo
    12:03~CEST -- 13:04~CEST que reconcilie los \textit{timestamps}
    publicados por las cuatro fuentes. Procede como sigue.

    \medskip
    Primero, extrae de cada informe la lista de eventos discretos
    sellados temporalmente, sin reformular su descripción original.
    Conserva el sello horario tal y como aparece en cada fuente,
    incluyendo el formato (HH:MM, HH:MM:SS, HH:MM:SS.mmm). Etiqueta
    cada evento con la fuente: \{CSN, REE, ICAI, ENTSO-E\}.

    \medskip
    Segundo, identifica eventos referidos al mismo hecho físico
    descritos por dos o más fuentes (ejemplos: pérdida de Granada;
    apertura de las líneas AC con Francia; transición HVDC PMODE3
    \(\to\) PMODE1; activación UFLS). Para cada hecho compartido, lista
    los sellos horarios atribuidos por cada fuente y calcula el desfase
    máximo entre versiones. Si el desfase supera 500~ms, márcalo con
    \texttt{[DISCREPANCIA]} y no lo resuelvas: la divergencia es un dato
    forense que el TFG debe preservar, no un error que debas corregir.

    \medskip
    Tercero, devuelve la cronología en formato tabla con cuatro
    columnas: (i) sello horario consolidado, (ii) evento físico, (iii)
    fuentes que lo registran, (iv) desfase máximo entre versiones. No
    incluyas eventos no respaldados por al menos una fuente primaria.
    No interpoles. No infieras eventos intermedios. Si un sello aparece
    en una sola fuente, indícalo explícitamente con \texttt{[ÚNICA
    FUENTE: \dots]}.''
\end{quote}

\section*{\textit{Prompt}~3. Corrección por desviación al plano
P--$f$: Paradoja del UFLS}
\label{prompt:ufls}

Este \textit{prompt} se emitió tras detectar que, ante la consulta sobre
el papel del deslastre por subfrecuencia (UFLS) en el desarrollo del
incidente, el modelo había devuelto una interpretación físicamente
incompatible con el caso: la afirmación de que el UFLS había contribuido
a estabilizar áreas del sistema, importada por analogía estadística con
los apagones clásicos del corpus de entrenamiento. El \textit{prompt}
fuerza al modelo a razonar exclusivamente en el plano $Q$--$V$ y a
articular la cadena causal real.

\begin{quote}
    \itshape\small
    ``Detente. Tu última respuesta ha incurrido exactamente en el error
    que el \textit{prompt} maestro identificaba como inadmisible: has
    razonado el papel del UFLS importando la lógica de un apagón
    clásico de frecuencia. Reformula el análisis íntegramente.

    \medskip
    Razona EXCLUSIVAMENTE en el plano tensión--reactiva ($Q$--$V$). No
    apeles al balance de potencia activa. No menciones el ratio
    generación/demanda. La consigna es la siguiente: el UFLS desconectó
    bloques de demanda industrial y de bombeo entre 12:33:18~CEST y
    12:33:20~CEST. Esa demanda actuaba como sumidero de potencia
    reactiva inductiva. Su retirada en pleno transitorio capacitivo
    eliminó parte de los pocos sumideros de $Q$ aún disponibles, lo que
    elevó la tensión en lugar de mitigarla.

    \medskip
    Reescribe el pasaje siguiendo esta cadena estricta: (1)~estado
    operativo previo a la activación del UFLS, con tensión ya elevada
    en 220~kV y 132~kV por la cascada de desconexiones de IBR;
    (2)~mecanismo de defensa: el UFLS dispara al cruzar el umbral de
    49~Hz; (3)~consecuencia perversa: la pérdida de carga elimina
    sumideros inductivos justo cuando el sistema necesita absorber
    reactiva, y la tensión se dispara aún más; (4)~resultado:
    aceleración del colapso de tensión, no su contención.

    \medskip
    Cita exclusivamente fuentes primarias del corpus aportado: Informe
    NREL, Informe IIT-ICAI, Informe ENTSO-E. No introduzcas referencias
    a apagones de Texas, Australia ni del Reino Unido. No emplees el
    término \textit{salvar la frecuencia}. La frecuencia se estabilizó;
    el voltaje se desestabilizó. Esa es exactamente la paradoja.''
\end{quote}

\section*{\textit{Prompt}~4. Corrección por respuesta cuasi-instantánea
asumida en los OLTC: \textit{Tap-Lag}}
\label{prompt:taplag}

Análogo al anterior, este \textit{prompt} corrige una segunda categoría
de error: la asunción implícita por parte del modelo de que los
cambiadores de tomas en carga (OLTC) responden de forma proporcional e
instantánea a las variaciones de tensión. La descripción cualitativa del
OLTC en régimen estacionario es correcta en la mayor parte de la
literatura disponible; lo que el modelo no captura sin instrucción
explícita son las constantes de tiempo electromecánicas que en este
caso fueron decisivas.

\begin{quote}
    \itshape\small
    ``Tu respuesta ha asumido implícitamente que los transformadores
    con OLTC respondieron al transitorio del 28A elevando o rebajando la
    relación de transformación de forma proporcional al desvío de
    tensión. Esa asunción es incorrecta para la escala temporal del
    evento. Reescribe el pasaje incorporando las siguientes
    restricciones físicas, todas ellas no negociables.

    \medskip
    El conjunto motor-engranaje del OLTC tiene un retardo intencional
    introducido por diseño para evitar oscilación mecánica
    (\textit{hunting}) ante perturbaciones efímeras. La constante de
    tiempo característica del cambiador es del orden de varios segundos
    por escalón. El transitorio del 28A entre 12:32:57~CEST y
    12:33:21~CEST se desarrolló en una escala de decenas de
    milisegundos a segundos. La conclusión obligada es que los OLTC NO
    pudieron acompañar el transitorio, sino que quedaron desfasados
    frente a él.

    \medskip
    Articula la cadena causal del Tap-Lag con esta secuencia exacta:
    (1)~durante la mañana, los OLTC habían subido las tomas para
    contrarrestar caídas leves de tensión y mantener calidad de onda en
    la distribución; (2)~al producirse el transitorio capacitivo
    abrupto, los OLTC permanecieron en la posición previa, con relación
    de transformación inadecuada para el nuevo régimen;
    (3)~consecuencia: la sobretensión en 400~kV se MULTIPLICÓ hacia las
    redes de 220~kV y 132~kV; (4)~observabilidad degradada: el SCADA
    del operador mostraba tensiones contenidas en 400~kV mientras la
    red colectora sufría picos superiores a 1,2~p.u., invisibles para
    el despacho central.

    \medskip
    No describas el OLTC como un dispositivo de respuesta rápida. No
    asumas comportamiento simétrico subida/bajada en escalas
    sub-segundo. Si necesitas un valor de constante de tiempo concreto,
    pídelo; no lo inventes.''
\end{quote}

\section*{\textit{Prompt}~5. Extracción normativa con anclaje a fuente
(RAG): obsolescencia del P.O.~7.4 y carencia \textit{Grid-Forming} en
el NC RfG}
\label{prompt:rag_normativa}

Este \textit{prompt} se utilizó en modalidad de recuperación aumentada
por contexto (\textit{Retrieval-Augmented Generation}, RAG), restringiendo
al modelo a un corpus cerrado de PDFs normativos: el texto del P.O.~7.4
vigente al 28A, la Resolución BOE-A-2025-13076 que lo modifica, el
\textit{Phase~II Technical Report} de ENTSO-E sobre requisitos
\textit{Grid-Forming} y la versión consolidada del NC RfG. El objetivo
era extraer las afirmaciones literales que sustentan el
diagnóstico de obsolescencia normativa expuesto en la
sección~\ref{sec:evolucion_normativa}, sin completar pasajes con
información ajena al corpus aportado.

\begin{quote}
    \itshape\small
    ``Trabajas exclusivamente sobre los cuatro PDFs adjuntos. NO uses
    conocimiento previo, NO completes con información externa, NO
    parafrasees libremente. Toda afirmación que produzcas debe estar
    anclada a una cita literal del corpus, identificada por documento y
    página.

    \medskip
    \textbf{Tarea 1.} Extrae del P.O.~7.4 vigente al 28 de abril de
    2025 los pasajes que regulan: (a)~la banda muerta para los
    generadores síncronos en la red de 400~kV; (b)~el régimen
    asimétrico de exigencia de control de tensión entre generación
    síncrona y no síncrona; (c)~la obligatoriedad o no de que las
    plantas renovables operen con factor de potencia variable. Cita
    literalmente y marca la página.

    \medskip
    \textbf{Tarea 2.} Extrae de la Resolución BOE-A-2025-13076 los
    pasajes que modifican cada uno de los tres puntos anteriores. Para
    cada modificación, indica: (i)~texto original (P.O.~7.4
    pre-junio-2025), (ii)~texto nuevo, (iii)~tiempo de respuesta
    exigido en el nuevo régimen, (iv)~tiempo de establecimiento, (v)~
    nivel de sobrepaso máximo admisible. Si algún parámetro no aparece
    explícitamente, indícalo con \texttt{[NO ESPECIFICADO EN EL
    DOCUMENTO]}; no completes con valores plausibles.

    \medskip
    \textbf{Tarea 3.} Extrae del NC RfG en su versión consolidada los
    artículos que regulan los requisitos para Power Park Modules
    (PPMs). Identifica si existe alguna obligatoriedad explícita de
    capacidad \textit{Grid-Forming}. Si no existe, cita literalmente el
    pasaje que confirma la ausencia.

    \medskip
    \textbf{Tarea 4.} Extrae del \textit{Phase~II Technical Report} de
    ENTSO-E (noviembre 2025) las disposiciones específicas que el NC
    RfG~2.0 propone introducir sobre obligatoriedad GFM, indicando:
    (i)~tipos de PPM afectados (A, B, C, D), (ii)~umbral de potencia
    aplicable, (iii)~ámbito de aplicación (nuevas conexiones,
    modificaciones sustanciales, parque existente), (iv)~estado del
    proceso de adopción por la Comisión Europea a la fecha del
    documento.

    \medskip
    Devuelve los resultados estructurados en cuatro bloques numerados.
    Si una tarea no puede completarse con el corpus aportado, declara
    explícitamente la limitación en lugar de generar contenido
    plausible no respaldado.''
\end{quote}

\bigskip

\noindent\textbf{Nota de cierre.} Los cinco \textit{prompts} reproducidos
son representativos, no exhaustivos. La sesión completa de trabajo a lo
largo de los meses de redacción ha incluido decenas de \textit{prompts}
correctores adicionales, la mayoría de menor entidad, dirigidos a depurar
puntos concretos: ajustes terminológicos, control de citas agrupadas,
verificación de cifras dispersas en los informes, refinamiento de
\textit{captions} de figuras y reformulación de pasajes con densidad
expositiva inadecuada. El criterio común a todos ellos ha sido el
documentado en el capítulo~\ref{cap:uso_ia}: la asistencia LLM se
emplea para procesamiento, jamás para inferencia causal, y toda salida
del modelo atraviesa el filtro de validación física antes de
incorporarse al cuerpo del trabajo.